\documentclass{zjureport}
% \special{dvipdfmx:config z 0} % 最终提交前可注释该行以恢复 PDF 压缩

% =============================================
% 报告信息：提交前请补全姓名、学号、教师等字段
% =============================================

\major{电子科学与技术}
\name{金大钧}
\title{本科实验报告}
\stuid{3230103274}
\college{信息与电子工程学院}
\date{}
\course{人工智能}
\instructor{金日成}
\expname{基于 Q-learning 的 Flappy Bird 智能体实验}
\renewcommand{\figurename}{图}
\renewcommand{\tablename}{表}
\renewcommand{\refname}{参考文献}

\begin{document}

\makecover
\makecontent
% \makeheader

\section{实验介绍}

\subsection{实验背景}

Flappy Bird 是一个规则简单但控制精度要求较高的游戏。玩家需要在连续移动的管道之间控制小鸟的高度，动作只有“不拍翅”和“拍翅”两种，但每一次动作都会影响后续一段时间内的小鸟速度和位置。因此，这一任务可以自然地看作一个序列决策问题：智能体在每一帧根据当前观测状态选择动作，并通过不断试错学习如何获得更高的游戏分数。

本实验采用说明文档给出的表格式 Q-learning 框架训练 Flappy Bird 智能体。智能体用字典保存状态--动作对的 $Q$ 值，并根据与环境交互得到的样本逐步更新 Q 表。实验效果与状态表示、训练参数和奖励设计密切相关。

\subsection{实验目标}

本实验的目标是在给定框架上实现 Q-learning 智能体，并通过实验比较不同设置下的模型表现：

\begin{enumerate}
    \item 补全 Q 表读取、未来回报估计、Q-learning 更新、动作选择和状态转换；
    \item 比较不同状态表示和 \texttt{obs\_mul\_factor} 设置；
    \item 调整 $\alpha$、$\epsilon$、$\gamma$ 及衰减策略；
    \item 尝试修改死亡惩罚和中心奖励；
    \item 保存表现最好的模型作为 \texttt{q\_best.pkl}。
\end{enumerate}

实验中同时记录参数配置、评估曲线和随机种子，用于分析训练过程中的波动。

\subsection{实验整体思路}

实验首先完成基础 Q-learning 训练流程；随后比较 \texttt{process\_obs()} 中的状态定义；之后调整学习率、探索率、折扣因子和衰减策略；最后尝试 reward 修改，并结合多种子结果分析训练波动。

由于表格式 Q-learning 直接以离散状态作为 Q 表索引，状态粒度会影响样本覆盖和 Q 值估计。因此，状态设计和参数调整需要结合实验结果一起分析。

\clearpage

\section{Q-learning 算法与 GameAI 实现}

\subsection{任务建模与 Q 表}

在本实验中，环境每一步返回观测值 \texttt{obs}，\texttt{process\_obs()} 将其转换为离散状态 $s$。智能体根据当前状态选择动作 $a$，环境返回下一观测、奖励 $r$ 和终止信息。动作空间为：

\begin{itemize}
    \item $a=0$：不拍翅；
    \item $a=1$：拍翅。
\end{itemize}

\texttt{GameAI} 使用字典 \texttt{self.q} 保存 Q 表，键为 \texttt{(state, action)}，值为对应的 $Q(s,a)$。状态由整数元组表示，未出现过的状态--动作对默认 $Q$ 值为 0。该逻辑由 \texttt{get\_q\_value()} 完成。

\subsection{Q-learning 更新与动作选择}

Q-learning 的核心思想是用当前样本修正状态--动作价值估计。对于一次交互得到的 $(s,a,r,s')$，更新公式为：

\[
Q(s,a) \leftarrow (1-\alpha)Q(s,a)+\alpha \left(r+\gamma \max_{a'}Q(s',a')\right)
\]

其中，$\alpha$ 是学习率，控制新样本对旧估计的影响；$\gamma$ 是折扣因子，控制未来奖励在当前决策中的重要程度；$\max_{a'}Q(s',a')$ 表示下一状态下能够取得的最大估计价值。

代码中，\texttt{best\_future\_reward()} 计算 $\max_{a'}Q(s',a')$，\texttt{update()} 按公式更新 Q 表。动作选择由 \texttt{choose\_action()} 完成：训练时使用 $\epsilon$-greedy 策略，在随机探索和选择当前最大 Q 值动作之间切换；评估时关闭随机探索。

\subsection{训练、评估与模型保存}

训练函数 \texttt{train()} 按照指定轮数反复运行游戏。每一局中，智能体根据当前状态选择动作，环境返回下一状态和奖励后，程序调用 \texttt{update()} 更新 Q 表。当游戏结束时进入下一局训练。

训练过程中可以按固定间隔调用 \texttt{evaluate()} 进行评估。训练结束后，\texttt{save\_q()} 将学习得到的 Q 表保存为模型文件，最终选取得分最高的模型作为 \texttt{q\_best.pkl}。

至此，代码形成了完整的 Q-learning 训练流程：观测处理、动作选择、环境交互、Q 表更新、评估和模型保存。后续实验均基于这一流程展开。

\clearpage

\section{Bonus1：状态表示设计与 obs 优化}

\subsection{示例状态与初始实现}

说明文档给出的示例方法是从 12 维观测值中选取目标管道位置、小鸟位置和小鸟速度，将状态表示为

\[
s=(x_{\mathrm{to\_pipe}},\,y_{\mathrm{to\_bottom}},\,v_y)
\]

其中 $x_{\mathrm{to\_pipe}}$ 表示小鸟最左端到最近管道最左端的水平距离，$y_{\mathrm{to\_bottom}}$ 表示小鸟顶端到下方管道边界的竖直距离，$v_y$ 表示竖直速度。连续观测值经 \texttt{obs\_mul\_factor} 放大取整后作为 Q 表键。

实验先按该示例思路实现三维状态，随后修改水平距离和目标管道切换方式。改进的动机是：能否通过管道主要取决于小鸟与管道边界的相对位置，左端到左端距离对接近管道时的几何关系刻画不足。

\subsection{水平与竖直距离定义实验}

早期实验固定 $\alpha=0.7,\gamma=0.95,\epsilon=0$，训练 50000 局，每 5000 局评估一次，最终评估为 10 局平均分。表 \ref{tab:obs-change} 列出各组实验内容。

\begin{table}[H]
    \centering
    \small
    \caption{水平与竖直距离定义实验}
    \label{tab:obs-change}
    \begin{tabular}{p{0.05\linewidth}p{0.45\linewidth}p{0.30\linewidth}r}
        \toprule
        编号 & 实验内容 & 主要变化 & 得分 \\
        \midrule
        O0 & 按说明文档示例实现三维状态：水平距离、下管道竖直距离、竖直速度 & 小鸟左端到管道左端 & 13.00 \\
        O1 & 在 O0 基础上修改水平距离和目标管道切换方式 & 改为基于边界位置计算 & 30.26 \\
        O2 & 在 O1 基础上修改竖直距离定义 & 小鸟底端到下管道边界 & 19.50 \\
        O3 & 在 O2 基础上增加到上管道边界的距离 & 三维状态变为四维状态 & 14.10 \\
        \bottomrule
    \end{tabular}
\end{table}

O1 得分最高，说明水平距离与目标管道切换方式对状态质量影响较大。相较 O0，O1 更接近实际通过管道时的边界关系，因此离散状态更容易对应到相近的动作选择。O2 将竖直距离改为小鸟底端到下管道边界，主要改变分箱边界，提升有限。O3 增加上管道距离后状态维度变为四维，但上下管道距离高度相关，反而增加状态稀疏性。

\subsection{在参数优化后重新考察离散化粒度}

早期 obs 修改后，实验先转向第四章的参数优化。原因是模型尚未稳定收敛时，继续增加观测量容易把状态设计和训练参数混在一起。确定较合理的训练主线后，再回到 \texttt{obs\_mul\_factor}。

\texttt{obs\_mul\_factor} 控制连续状态的离散化粒度。50k 基础组统一使用 $\alpha=0.5$、count 学习率衰减 $p=0.4$、$\gamma=0.97$、初始 $\epsilon=0.2$、最小 $\epsilon=0.01$、指数探索衰减时间常数 10000；优化组将训练轮数增至 100k，并将探索衰减时间常数放慢到 20000。

\begin{table}[H]
    \centering
    \small
    \caption{离散化粒度实验：同一状态结构下改变 \texttt{obs\_mul\_factor}}
    \label{tab:obs-factor}
    \begin{tabular}{p{0.08\linewidth}rrrrr}
        \toprule
        编号 & $k$ & 训练轮数 & $\epsilon$ 衰减常数 & 最终均分 & 种子范围 \\
        \midrule
        F1 & 10 & 50k & 10000 & 38.18 & 6.67--83.09 \\
        F2 & 20 & 50k & 10000 & 110.04 & 73.12--167.29 \\
        F3 & 30 & 50k & 10000 & 256.07 & 125.42--480.91 \\
        F4 & 40 & 50k & 10000 & 115.04 & 53.24--220.43 \\
        F5 & 50 & 50k & 10000 & 39.63 & 23.07--54.86 \\
        F6 & 40 & 100k & 20000 & 750.82 & 210.24--1533.38 \\
        F7 & 50 & 100k & 20000 & 79.96 & 65.89--102.56 \\
        \bottomrule
    \end{tabular}
\end{table}

50k 条件下，$k=30$ 最好；$k=40,50$ 反而下降，说明状态划分过细会使访问次数不足。F6 表明 $k=40$ 在训练轮数和探索时间增加后可以发挥优势；F7 则说明 $k=50$ 的状态空间过大，100k 训练仍不足以覆盖。最终采用 \texttt{obs\_mul\_factor=40}，但它依赖更长训练、较慢探索衰减和 count 学习率衰减，并非单独修改 factor 即可生效。

\clearpage

\section{参数调优实验}

状态表示调整后，继续比较学习率 $\alpha$、探索率 $\epsilon$、探索率衰减、折扣因子 $\gamma$ 和学习率衰减。实验先用固定 $\alpha/\epsilon$ 确定可用范围，再引入衰减机制处理前期探索和后期稳定之间的矛盾。

\subsection{固定学习率与固定探索率}

第一组实验固定 $\gamma=0.95$，训练 50000 局，每 5000 局评估一次，每次评估 100 局。粗调先在较大的 $\alpha/\epsilon$ 网格中搜索，结果见表 \ref{tab:alpha-epsilon-coarse}。

\begin{table}[H]
    \centering
    \small
    \caption{固定 $\alpha$ 与固定 $\epsilon$ 的粗调结果}
    \label{tab:alpha-epsilon-coarse}
    \begin{tabular}{crrrr}
        \toprule
        $\alpha \backslash \epsilon$ & 0 & 0.02 & 0.05 & 0.10 \\
        \midrule
        0.10 & 46.14 & 34.74 & 25.22 & 9.48 \\
        0.20 & 34.12 & 34.26 & 76.99 & 28.46 \\
        0.30 & 10.10 & 39.73 & 24.03 & 9.78 \\
        0.50 & 70.88 & 9.85 & 12.73 & 5.41 \\
        0.70 & 26.80 & 5.95 & 2.67 & 1.69 \\
        \bottomrule
    \end{tabular}
\end{table}

粗调结果显示，较大的探索率通常降低最终评估分，$\epsilon=0.10$ 在多数学习率下表现较差；较大的学习率也不稳定，$\alpha=0.7$ 基本失效。随后在较小学习率和低探索率附近细调，结果见表 \ref{tab:alpha-epsilon-fine}。

\begin{table}[H]
    \centering
    \small
    \caption{固定 $\alpha$ 与固定 $\epsilon$ 的细调结果}
    \label{tab:alpha-epsilon-fine}
    \begin{tabular}{crrr}
        \toprule
        $\alpha \backslash \epsilon$ & 0.01 & 0.02 & 0.03 \\
        \midrule
        0.05 & 27.91 & 13.49 & 10.68 \\
        0.08 & 64.40 & 112.06 & 47.36 \\
        0.10 & 49.86 & 34.74 & 42.47 \\
        0.12 & 41.32 & 92.89 & 49.02 \\
        0.15 & 36.75 & 26.97 & 63.94 \\
        \bottomrule
    \end{tabular}
\end{table}

细调中 $\alpha=0.08,\epsilon=0.02$ 得分最高，$\alpha=0.12,\epsilon=0.02$ 次之。两轮固定参数实验只能确定可用范围：低探索率有助于补充状态覆盖，但固定 $\epsilon$ 不能同时满足前期探索和后期稳定，因此后续继续测试探索率衰减。

\subsection{探索率衰减}

第二组实验固定 $\alpha=0.1,\gamma=0.95$，训练 50000 局，先比较固定探索率、线性衰减、乘数衰减和指数衰减。表中“后期均分”为最后三次周期评估的平均值，用于观察训练后段是否仍有明显震荡。

探索率衰减控制训练第 $t$ 局使用的 $\epsilon_t$。线性衰减为
\[
\epsilon_t=\epsilon_{\min}+(\epsilon_0-\epsilon_{\min})\max(0,1-t/T),
\]
指数衰减为
\[
\epsilon_t=\epsilon_{\min}+(\epsilon_0-\epsilon_{\min})e^{-t/\tau},
\]
乘数衰减为
\[
\epsilon_t=\max(\epsilon_{\min},\epsilon_0\lambda^t).
\]
其中 $\epsilon_{\min}$ 是最小探索率，$T$ 为线性衰减轮数，$\tau$ 为指数衰减时间常数，$\lambda$ 为乘数衰减系数。

\begin{table}[H]
    \centering
    \small
    \caption{探索率衰减方式比较}
    \label{tab:epsilon-decay-method}
    \begin{tabular}{p{0.08\linewidth}p{0.48\linewidth}rr}
        \toprule
        编号 & 探索率设置 & 最终得分 & 后期均分 \\
        \midrule
        E1 & 固定 $\epsilon=0.02$ & 34.74 & 34.46 \\
        E2 & 固定 $\epsilon=0.10$ & 9.48 & 7.12 \\
        E3 & 线性衰减：$0.10\rightarrow0.01$，30000 局衰减 & 55.25 & 93.10 \\
        E4 & 乘数衰减：初始 0.10，系数 0.9998，下限 0.01 & 52.47 & 42.41 \\
        E5 & 乘数衰减：初始 0.10，系数 0.99992，下限 0.01 & 63.88 & 45.97 \\
        E6 & 指数衰减：初始 0.10，时间常数 5000，下限 0.01 & 13.94 & 27.17 \\
        E7 & 指数衰减：初始 0.10，时间常数 10000，下限 0.01 & 55.55 & 130.87 \\
        \bottomrule
    \end{tabular}
\end{table}

固定大探索率 E2 后期表现较差，说明随机动作过多会破坏已学到的策略。指数衰减的两个初始候选中，时间常数 5000 的 E6 衰减过快，最终得分和后期均分都较低；时间常数 10000 的 E7 后期均分最高。因此后续选择指数衰减作为主线，并继续围绕初始探索率、最小探索率和时间常数做细化测试，结果见表 \ref{tab:epsilon-exp-tune}。

\begin{table}[H]
    \centering
    \small
    \caption{指数探索率衰减细化实验}
    \label{tab:epsilon-exp-tune}
    \begin{tabular}{p{0.08\linewidth}p{0.46\linewidth}rr}
        \toprule
        编号 & 指数衰减设置 & 最终得分 & 后期均分 \\
        \midrule
        E8 & 初始 0.10，时间常数 7500，下限 0.01 & 48.59 & 50.58 \\
        E9 & 初始 0.10，时间常数 12500，下限 0.01 & 38.22 & 38.82 \\
        E10 & 初始 0.10，时间常数 15000，下限 0.01 & 37.11 & 34.70 \\
        E11 & 初始 0.10，时间常数 10000，下限 0.005 & 54.91 & 47.37 \\
        E12 & 初始 0.15，时间常数 10000，下限 0.01 & 52.78 & 34.80 \\
        E13 & 初始 0.20，时间常数 10000，下限 0.01 & 48.09 & 35.16 \\
        \bottomrule
    \end{tabular}
\end{table}

细化实验没有给出单调结论，但可以排除过快衰减。后续主线采用 E13 的设置，即初始 $\epsilon=0.2$、时间常数 10000、最小 $\epsilon=0.01$。在状态空间继续增大或训练轮数变化时，再根据访问不足或后期随机动作过多的问题调整衰减速度。

\subsection{折扣因子 $\gamma$}

第三组实验开始使用随机种子 42、100 和 2026 比较稳定性。其他参数为 $\alpha=0.1$、初始 $\epsilon=0.2$、最小 $\epsilon=0.01$、指数探索衰减、时间常数 10000、训练 50000 局。表中“最终均分”为三个种子最终得分的平均，“种子范围”为三个最终得分的最小值和最大值；多种子“后期均分”为各个种子最后三次周期评估均值再取平均。

\begin{table}[H]
    \centering
    \small
    \caption{折扣因子实验}
    \label{tab:gamma}
    \begin{tabular}{p{0.08\linewidth}rrrr}
        \toprule
        编号 & $\gamma$ & 最终均分 & 种子范围 & 后期均分 \\
        \midrule
        G1 & 0.80 & 9.24 & 5.06--15.25 & 13.22 \\
        G2 & 0.90 & 51.09 & 16.38--77.80 & 34.54 \\
        G3 & 0.95 & 718.81 & 36.36--2052.40 & 294.10 \\
        G4 & 0.97 & 113.42 & 27.12--179.07 & 93.15 \\
        G5 & 0.98 & 73.42 & 41.51--91.15 & 76.11 \\
        G6 & 0.99 & 70.74 & 21.29--111.11 & 62.03 \\
        G7 & 1.00 & 21.88 & 13.49--28.58 & 19.66 \\
        \bottomrule
    \end{tabular}
\end{table}

$\gamma$ 过低会削弱远期奖励传播；$\gamma=1$ 又容易放大远期估计误差。G3 均分最高但种子范围极大，主要由单个高分种子拉高。G4 的后期表现和稳定性更适合作为后续实验基础，因此选择 $\gamma=0.97$。

\subsection{学习率衰减}

第四组实验采用 $\gamma=0.97$、初始 $\epsilon=0.2$、最小 $\epsilon=0.01$、指数探索衰减时间常数 10000，比较固定学习率、全局衰减和按状态--动作访问次数衰减。

学习率衰减控制 TD 更新中的步长。全局线性衰减和指数衰减分别为
\[
\alpha_t=\alpha_{\min}+(\alpha_0-\alpha_{\min})\max(0,1-t/T_\alpha),
\qquad
\alpha_t=\alpha_{\min}+(\alpha_0-\alpha_{\min})e^{-t/\tau_\alpha}.
\]
count 衰减不只依赖训练轮数，而是按状态--动作对的访问次数调整：
\[
\alpha(s,a)=\frac{1}{N(s,a)^p}.
\]
因此常见状态的更新步长会更快下降，低频状态仍保留较大的修正幅度。

\begin{table}[H]
    \centering
    \small
    \caption{学习率衰减实验}
    \label{tab:alpha-decay}
    \begin{tabular}{p{0.08\linewidth}p{0.38\linewidth}rrr}
        \toprule
        编号 & 学习率设置 & 最终均分 & 种子范围 & 后期均分 \\
        \midrule
        A1 & 固定 $\alpha=0.1$ & 113.42 & 27.12--179.07 & 93.15 \\
        A2 & 固定 $\alpha=0.5$ & 11.11 & 6.58--13.54 & 13.49 \\
        A3 & 线性衰减：$0.5\rightarrow0.01$ & 129.86 & 100.89--157.64 & 134.94 \\
        A4 & 指数衰减：$0.5\rightarrow0.01$，时间常数 10000 & 222.01 & 143.48--346.41 & 174.42 \\
        A5 & 指数衰减：$0.5\rightarrow0.05$，时间常数 10000 & 209.99 & 137.48--246.47 & 178.64 \\
        A6 & count 衰减：$\alpha=1/N(s,a)^{0.4}$ & 256.07 & 125.42--480.91 & 298.19 \\
        A7 & count 衰减：$\alpha=1/N(s,a)^{0.6}$ & 40.36 & 38.73--42.53 & 39.70 \\
        A8 & count 衰减：$\alpha=1/N(s,a)$ & 38.71 & 21.36--71.36 & 28.44 \\
        \bottomrule
    \end{tabular}
\end{table}

固定高学习率 A2 明显不稳定。全局衰减能减小后期震荡，但无法区分常见状态和罕见状态。A6 的 count 衰减效果最好：高频状态逐渐降低步长，低频状态保留较大更新幅度。A7、A8 衰减过快，许多状态尚未充分学习就接近冻结。后续主线采用 $\gamma=0.97$、初始 $\epsilon=0.2$、最小 $\epsilon=0.01$、指数探索衰减和 count 学习率衰减 $p=0.4$。

\clearpage

\section{Bonus2：Reward 设计与修改实验}

环境原始奖励包括：存活每帧 $+0.1$，通过管道 $+1.0$，死亡 $-1.0$，触碰上边界 $-0.5$。Reward 修改会直接改变 TD 更新中的即时回报项，因此应从分别考察死亡惩罚、直接中心奖励和基于势函数的中心奖励。

\subsection{死亡惩罚修改}

说明文档中的示例代码将死亡奖励从 $-1$ 放大为 $-1000$，目的是使智能体更重视避免死亡。实验首先比较不同死亡惩罚。该组固定 $\alpha=0.1,\gamma=0.95$，初始 $\epsilon=0.2$，最小 $\epsilon=0.01$，指数探索衰减时间常数 10000，训练 50000 局，使用 3 个训练种子。

\begin{table}[H]
    \centering
    \small
    \caption{死亡惩罚实验}
    \label{tab:death-penalty}
    \begin{tabular}{p{0.1\linewidth}rrrr}
        \toprule
        编号 & 死亡惩罚 & 最终均分 & 种子范围 & 后期均分 \\
        \midrule
        R1 & $-10$ & 45.04 & 25.91--63.39 & 43.15 \\
        R2 & $-100$ & 718.81 & 36.36--2052.40 & 294.10 \\
        R3 & $-500$ & 41.14 & 29.85--52.08 & 41.46 \\
        R4 & $-1000$ & 42.71 & 39.40--48.09 & 38.33 \\
        \bottomrule
    \end{tabular}
\end{table}

$-100$ 的平均分最高，但方差也最大，说明它更容易在某些训练轨迹中形成高分策略，并不代表所有种子都稳定收敛。过大的死亡惩罚会使失败样本在更新中占据过高权重，导致临界状态的 $Q$ 值过度悲观；过小惩罚又会削弱死亡信号。综合后续实验，最终代码中采用 $-100$。

\subsection{直接中心奖励}

死亡惩罚只处理失败状态，没有直接鼓励小鸟接近管道中心。为测试这一想法，实验加入中心位置奖励。设管道缺口中心为 $y_c$，小鸟中心为 $y_b$，安全半高度为 $h$，则中心分数近似为

\[
1-\min\left(\frac{|y_c-y_b|}{h},2\right)
\]

同时用水平距离权重 $w_x(s)$ 限制该奖励主要在接近目标管道时起作用。实际加入的势函数为

\[
\Phi(s)=c\,w_x(s)\left(1-\min\left(\frac{|y_c-y_b|}{h},2\right)\right),
\]

其中 $c$ 为中心奖励系数。直接中心奖励将该项直接加到环境奖励中：

\[
r'=r+\Phi(s')
\]

表 \ref{tab:center-direct} 中 D0--D3 使用相同训练设置：\texttt{obs\_mul\_factor=40}，训练 100k，$\gamma=0.97$，初始 $\epsilon=0.2$，最小 $\epsilon=0.01$，探索衰减时间常数 20000，count 学习率衰减 $p=0.4$。

\begin{table}[H]
    \centering
    \small
    \caption{直接中心奖励实验}
    \label{tab:center-direct}
    \begin{tabular}{p{0.1\linewidth}rrrr}
        \toprule
        编号 & 中心奖励系数 & 最终均分 & 种子范围 & 后期均分 \\
        \midrule
        D0 & 0 & 750.82 & 210.24--1533.38 & 525.45 \\
        D1 & 0.01 & 485.90 & 181.64--1043.39 & 381.97 \\
        D2 & 0.02 & 410.43 & 159.65--771.74 & 360.63 \\
        D3 & 0.05 & 195.73 & 128.23--311.17 & 218.21 \\
        \bottomrule
    \end{tabular}
\end{table}

直接中心奖励未表现出稳定提升。原因是“当前更接近中心”不等价于“到达管道时能安全通过”。Flappy Bird 的动作有惯性，当前高度、速度和距离共同决定未来位置；单独奖励当前居中可能使智能体偏向短期居中，而不一定学到正确的提前拍翅时机。

\subsection{基于势函数的中心奖励}

直接加奖励可能改变原任务的最优策略。为减少这一问题，又测试了基于势函数的奖励塑形：

\[
r'=r+\gamma\Phi(s')-\Phi(s)
\]

理论上，形如 $\gamma\Phi(s')-\Phi(s)$ 的塑形项可以改变学习过程而不改变最优策略。实验中使用同样的中心势函数，并在终止状态令下一状态势函数为 0。训练设置为 100k，$\gamma=0.97$，初始 $\epsilon=0.2$，最小 $\epsilon=0.01$，探索衰减时间常数 30000，count 学习率衰减 $p=0.4$。

\begin{table}[H]
    \centering
    \small
    \caption{基于势函数的中心奖励实验}
    \label{tab:center-potential}
    \begin{tabular}{p{0.1\linewidth}rrrr}
        \toprule
        编号 & 势函数系数 & 最终均分 & 种子范围 & 后期均分 \\
        \midrule
        P0 & 0 & 435.93 & 105.08--630.96 & 380.26 \\
        P1 & 0.005 & 332.35 & 174.05--472.08 & 282.89 \\
        P2 & 0.008 & 382.34 & 134.49--541.44 & 356.57 \\
        P3 & 0.01 & 216.97 & 181.90--238.25 & 191.42 \\
        P4 & 0.02 & 284.92 & 151.49--525.12 & 280.28 \\
        \bottomrule
    \end{tabular}
\end{table}

非零势函数系数均未超过 P0。一个重要原因是，本实验的势函数由原始连续观测计算，而 $Q$ 表状态是离散后的三维元组；势函数所表达的“接近中心”并不完全等价于 Q-learning 实际使用的状态价值结构。另外，中心距离没有显式考虑速度和到达管道所需时间，因此对拍翅时机的指导有限。最终模型不采用额外中心奖励。

\clearpage

\section{高分模型与训练随机性讨论}

\subsection{最高分模型与稳定性评价}

最终模型选择基于已经训练出的模型评估分，因此单个模型的最高分具有直接意义。实验中最高的单模型评估分来自较早阶段的 reward 实验：死亡惩罚设为 $-100$、$\gamma=0.95$ 时，一个训练种子的最终评估平均分达到 2052.40。该结果明显高于后续大多数优化实验，是最终提交模型的重要候选。

但方法分析不能只依据单次最高值。同一配置下另外两个训练种子的最终评估分只有 36.36 和 67.68，说明 2052.40 不是稳定复现的平均表现，而是一次成功训练轨迹形成的高分策略。因此，该模型适合作为最高分提交候选，但该参数组合本身不能仅由这一点说明稳定有效。

\subsection{训练随机性的来源}

同一参数下不同种子差异很大，主要来自训练早期的访问路径差异。Flappy Bird 的状态转移对动作序列敏感，早期一次随机拍翅会改变后续高度、速度和接近管道时的状态。表格式 Q-learning 又只更新实际访问到的状态--动作对，因此不同种子会产生不同的 Q 表覆盖。

如果某次早期探索进入较优轨迹，相关状态会被反复访问和强化，形成正反馈；如果早期轨迹较差，模型可能长期停留在低分策略附近。离散状态空间增大后，这种差异更明显，因为很多状态--动作对在有限训练轮数内访问次数不足。

\subsection{后期优化结果的解释}

后期优化的主要作用是提高多种子平均表现和训练稳定性，不一定重新产生早期的极端高分。例如 \texttt{obs\_mul\_factor=40} 配合 100k 训练后，多种子均分达到 750.82，明显高于多数中间实验，但仍低于 2052.40 的单次最高值。这说明后期主线改善了常见训练轨迹下的表现，但没有再次触发同等水平的高分轨迹。

第五章的 center reward 结果也印证这个角度理解。直接中心奖励实验中，$coef=0$ 的 baseline 均分较高，部分原因是其中一个种子达到 1533.38 分；非零中心奖励没有稳定超过该 baseline，但部分配置仍出现过较高分数。因此当前结论不是 center reward 完全无效，而是现有中心势函数未表现出稳定优势。它只基于当前位置，未充分考虑速度和到达管道时刻的位置，对拍翅时机的指导有限。

\subsection{可重复性措施}

为区分偶然高分和稳定改进，后期实验通常使用 100 局评估和 3 个训练种子。训练时设置 Python \texttt{random}、NumPy 和环境 seed；评估时使用固定的 \texttt{seed=42+i} 序列，并关闭 $\epsilon$ 随机探索，以减少评估阶段的随机性。

每组实验通过 \texttt{--exp-name} 保存到独立目录，\texttt{results.txt} 记录参数、随机种子、训练轮数、评估间隔和评估曲线。Git 用于记录代码修改，便于区分不同阶段的实现。上述处理不改变算法本身，但使高分模型来源、参数变化和多种子结果可以追踪。

\clearpage

\section{实验总结}

本实验完成了基于 Q-learning 的 Flappy Bird 智能体，实现了 Q 表更新、$\epsilon$-greedy 动作选择、训练评估流程和模型保存。在此基础上，实验进一步完成了状态表示设计、超参数调优和 reward 修改，并通过多种子实验分析训练稳定性。

实验表明，表格式 Q-learning 的效果高度依赖状态离散化和访问次数。过粗的状态表示会混淆不同局面，过细的状态表示又会导致访问稀疏；学习率、探索率和衰减方式需要与状态空间大小共同调整。最终主线采用更细的状态离散化、指数探索衰减和 count 学习率衰减，在多种子平均表现上优于多数中间实验。

Reward 修改方面，死亡惩罚对学习方向影响明显，最终采用 $-100$；直接中心奖励和基于势函数的中心奖励没有表现出稳定优势。实验也显示，单次最高分和多种子稳定性是不同评价维度：2052.40 的高分模型可作为提交候选，但方法优劣仍需要结合均值、种子范围和训练过程判断。

\clearpage

\end{document}
